\newcommand{\aphorism}[3]{%
  \vspace{3ex plus 1ex minus .2ex}%
  \noindent\parbox{.85\textwidth}{\raggedright\itshape #1 \\ 
    \hspace*{\fill}{\upshape -- #2, \MakeUppercase{#3}}}}

\newcommand{\header}[1]{\vspace{1ex}\par\noindent\textbf{#1}\hspace{0.618em plus 0.1em minus 0.05em}}


\chapter*{Preface}
The Theory of Computation is awfully nice.
It is beautiful, timely,
and has deep connections to other areas inside of computing,
to mathematics, and to the wider intellectual world.

The subject makes a delightful course.
Students learn both of the power of comptuation, and
of its limits.
There are underlying principles, but there are also as-yet
undiscovered areas within easy reach of a first course.
This text aims to mirror that, to be both rigorous and delightful.

We cover the standard curriculum, but in a way that is different
from standard texts, emphasizing intuition and background, as well
as arranging the material to make it possible to express to students
the current research frontier.




\header{For students}
The Theory of Computation has nearly a hundred year history.
The subject tells a story, from the initial settling on a precise definition,
through unsolvability results and a branching into areas such 
as automata, up to today's interest in complexity.

Here you will see that story, not in its full detail of course, but in
sufficient depth that you can understand both its main ideas and how
it applies to various other areas.

A liberal approach to education, according to a 
\citetext{favorite definition of mine}{\protect\url{https://www.aacu.org/leap/what-is-a-liberal-education}},
``empowers individuals and prepares them to deal with complexity, diversity, and change. This approach emphasizes broad knowledge \ldots{} as well as in-depth achievement in a specific field of interest.'' 
% https://www.aacu.org/leap/what-is-a-liberal-education
Here you will see the story told in a liberal way, but without stinting
on needed technical detail.




\header{For teachers}
This text covers the standard topics: Turing machines, unsolvability, 
automata, languages, and complexity.
But the approach of this book differs from the standard one 
in a number of ways,
all of which aim to bring out how awfully nice the subject is.

First, we start with Turing machines.
These are after all the touchstone.
We give a rigorous definition, but we also have an extensive discussion
of the motivation for that definition.
And, we also include an extensive discussion of Church's Thesis, 
taking care to include discussion of those point that experience with
questions on the Internet shows are points of frequent confusion. 

This approach illustrates the overall 
liberal program of not being shy about stopping, where stopping makes sense, 
to consider an interesting point, for its pure interest.

\citetext{Winklmann and Schauble}{\autocite{WinklmannSchauble}} 
use many familiar 
constructs of practical computing to illustrate the 
principles that we are studying.
Surely basing a development on the prior work students have done
and tying new ideas to
their existing knowledge are sound teaching principles.
\citetext{Meyer}{\autocite{Meyer}} uses Scheme as a model of computing, 
as we here use the Turing machine
or as another text might use a register machine.
Surely expressing our thoughts clearly and succinctly is also a 
sound teaching principle.

% We don't quite do either of these. 
We use Scheme as a single, uniform, practical language in which we 
\citetext{express our 
algorithmic thoughts}{John McCarthy, 
the inventor of Scheme's ancestor Lisp, suggested the language
as well-suited for the Theory of Computation.
Two arguments in that direction 
are the language's elegance and expresiveness.
This is captured  an online obituary of
McCarthy written by Bertrand Meyer for the 
\textit{Communications of the Association for Computing Machinery}:
``The Lisp 1.5 manual, published in 1962, was another masterpiece; 
as early as page 13 it introduces\Dash an unbelievable feat, especially 
considering that the program takes hardly more than half a page\Dash an 
interpreter for the language being defined, written in that very language! 
The more recent reader can only experience here the kind of visceral, 
poignant and inextinguishable jealously that overwhelms us the first time we 
realize that we will never be able to attend the premi\`ere of Don Giovanni at 
the Estates Theater in Prague on 29 October, 1787 \ldots\,. 
What may have been the reaction of someone in `Data Processing,' 
such as it was in 1962, suddenly coming across such a language manual?'' 
\autocite{MeyerB}}.


\aphorism{It is common for people first starting to grapple with
computers to make large-scale computations of things they might have done on a
smaller scale by hand. They might print out a table of the first
10,000 primes,
only to find that their printout isn't something they really wanted after all. 
They discover by this kind of experience that what they really want is 
usually not some collection of ``answers''\Dash what they want is
understanding.}{W Thurston}{ON PROOF AND PROGRESS IN MATHEMATICS}

\vspace{1ex plus .25ex minus .1ex}
\aphorism{[S]cience has made nature intelligible in terms of its symmetries
and regularities, analyzing its most wayward forms into components of
a regular and measurable shape.
\ldots
As a result we tend to see nature and to deal with it as an ``order''
from which the element of spontaniety has been``screened out.''
But this order is \textit{maya}, and the ``true suchness'' of things
has nothing in common with the purely conceptual aridities of perfect squares, 
circles, or triangles\Dash except by spontaneous accident.
Yet this is why the \ldots\ mind is dismayed when ordered conceptions 
of the universe break down, and when the basic behavior of the physicial
world is found to the a ``principle of uncertainty.''}{A Watts}{THE WAY OF ZEN}
% from p 180

\vspace{1ex plus .25ex minus .1ex}
\aphorism{Teach concepts, not tricks}{Gian-Carlo Rota}{TEN LESSONS I WISH I HAD LEARNED BEFORE I STARTED TEACHING
DIFFERENTIAL EQUATIONS}
% from an MAA talk

\vspace{2ex plus 1\fill}
\begin{flushright}
  \begin{tabular}{l@{}}
    Jim Hef{}feron \\ 
    Saint Michael's College \\
    Colchester, VT USA \\
    \texttt{joshua.smcvt.edu/computing} \\
    \ymd{2017}{01}{01}
  \end{tabular}
\end{flushright}
